\documentclass[11pt]{article}
\usepackage[a4paper,margin=1in]{geometry}
\usepackage{amsmath,amssymb,amsthm,booktabs,enumitem,mathtools,microtype}
\usepackage[T1]{fontenc}
\usepackage{lmodern}
\usepackage[hidelinks]{hyperref}

\newtheorem{theorem}{Theorem}[section]
\newtheorem{proposition}[theorem]{Proposition}
\newtheorem{lemma}[theorem]{Lemma}
\newtheorem{corollary}[theorem]{Corollary}
\theoremstyle{definition}
\newtheorem{definition}[theorem]{Definition}
\theoremstyle{remark}
\newtheorem{remark}[theorem]{Remark}

\newcommand{\mob}{\mu}
\newcommand{\TT}{\{-1,0,1\}}
\newcommand{\Fcal}{\mathcal F}
\newcommand{\e}{\mathrm e}
\newcommand{\dd}{\delta}
\newcommand{\thh}{\vartheta}
\newcommand{\abs}[1]{\lvert#1\rvert}
\newcommand{\norm}[1]{\lVert#1\rVert}

\hypersetup{
  pdftitle={Polylogarithmic-Clock Uniformization for Phasewise Chowla-Free Lag-Two Memory},
  pdfauthor={RH research program},
  pdfsubject={A uniform Mobius limit for polylogarithmically growing finite clocks},
  pdfkeywords={Mobius function, Davenport estimate, growing clocks, lag-two memory}
}

\title{Polylogarithmic-Clock Uniformization for\\
Phasewise Chowla-Free Lag-Two Memory}
\author{RH research program}
\date{August 8, 2026}

\begin{document}
\maketitle

\begin{abstract}
Fix $B>0$.  For every clock
$1\le q\le\lfloor(\log N)^B\rfloor$, consider the RH-379 class of
universally distance-two-safe, $q$-periodic lag-two sign tables whose
shift-two interpolation coefficient vanishes separately at each phase.
We prove a limit uniform in both the clock and every table family in that
class.  If $S_N(q,f)$ is its normalized M\"obius score and $L_q(f)$ is the
fixed-clock limit, then
\[
 \sup_{q\le\lfloor(\log N)^B\rfloor}\ \sup_{f\in\Fcal_q}
       \abs{S_N(q,f)-L_q(f)}\longrightarrow0.
\]
The proof truncates squarefreeness at
$P=\lfloor\sqrt{\log\log N}\rfloor$.  With
$M_P=(\prod_{p\le P}p)^2$, $Q=\operatorname{lcm}(q,M_P)$,
$\tau_P=\sum_{p>P}p^{-2}$, and $D_*(N)$ the maximum of the uniform
Davenport sums at $N$ and $N-2$, we retain the explicit bound
\[
 \abs{S_N(q,f)-L_q(f)}
 \le \frac{4\sqrt QD_*(N)}{N}+13\tau_P+\frac{6Q}{N}+\frac4N.
\]
The Fourier coefficient costs are $1,1,2$; the last factor cannot be
removed because a legal self-compatible table has $c_{21}=-2$.
Consequently the finite optimizers converge uniformly to the fixed-clock
optimizers, their triangular-array maximum tends to the RH-379 endpoint
$B_\infty$, and the cofinal square clocks supply an explicit positive
diagonal witness after the admissible set becomes nonempty.  The result is
not polynomial-clock uniformity, an adaptive-capacity limit, a projective
infinite selector, an operator or trace construction, or a statement about
Riemann zeros.
\end{abstract}

\noindent\textbf{Keywords:} M\"obius function; Davenport estimate;
polylogarithmic clocks; squarefree cutoff; phasewise lag-two memory.

\section{Frozen problem and claim boundary}

RH-366 records Davenport's estimate uniformly in the additive frequency
\cite{Davenport1937,IwaniecKowalski2004,RH366}.  RH-378 gives the exact
six-monomial interpolation of lag-two tables, including the ordinary
shift-two coefficient $c_{11}$ \cite{RH378}.  RH-379 then optimizes the
universally safe class with $c_{11}(r)=0$ at every fixed phase and proves
\begin{equation}
                         \sup_{q<\infty}G(q)=B_\infty.
 \label{eq:rh379-supremum}
\end{equation}
That theorem fixes $q$ before $N\to\infty$ \cite{RH379}.  The purpose here
is exactly one limit exchange: clocks may now grow no faster than a fixed
power of $\log N$.

Let
\begin{equation}
 \mob_0(m)=
 \begin{cases}
  \mob(m),&m\ge1,\\
  0,&m\le0.
 \end{cases}
 \label{eq:zero-padding}
\end{equation}
The zero padding in \eqref{eq:zero-padding} is part of the finite score,
not an asymptotic convention.

\begin{definition}[Phasewise class]
Fix $q\ge1$.  A family $f=(f_r)_{r\bmod q}$ belongs to $\Fcal_q$ when
each $f_r:\TT^2\to\{-1,+1\}$ and the following hold.
\begin{enumerate}[leftmargin=2em]
 \item For every ternary input word, the output
 $\epsilon_n=f_{n\bmod q}(a_{n-2},a_n)$ never has
 $\epsilon_n=\epsilon_{n+2}=+1$.
 \item In the unique interpolation
 \begin{equation}
 zf_r(x,z)=c_{01}(r)z+c_{02}(r)z^2+c_{11}(r)xz
 +c_{12}(r)xz^2+c_{21}(r)x^2z+c_{22}(r)x^2z^2,
 \label{eq:interpolation}
 \end{equation}
 one has $c_{11}(r)=0$ for every phase $r$.
\end{enumerate}
\end{definition}

The first condition is universal safety, not safety merely along the
M\"obius word.  The second is phasewise: no cancellation between unknown
shift-two correlations carrying different phase weights is assumed.

Define
\begin{equation}
 S_N(q,f)=\frac1N\sum_{n\le N}\mob(n)
 f_{n\bmod q}\bigl(\mob_0(n-2),\mob(n)\bigr).
 \label{eq:finite-score}
\end{equation}
The one-site progression densities come from RH-375, and RH-379 supplies
the lag-two squarefree density and the resulting fixed-clock limit
\cite{RH375,RH379}:
\begin{align}
 \dd_{q,r}&=\lim_{N\to\infty}\frac1N
  \sum_{\substack{n\le N\\n\equiv r\ (q)}}\mob(n)^2,
 \label{eq:delta}\\
 \thh_{q,r}&=\lim_{N\to\infty}\frac1N
  \sum_{\substack{n\le N\\n\equiv r\ (q)}}
       \mob(n-2)^2\mob(n)^2.
 \label{eq:theta}
\end{align}
The corresponding fixed-clock functional is
\begin{equation}
 L_q(f)=\sum_{r\bmod q}
       \bigl(c_{02}(r)\dd_{q,r}+c_{22}(r)\thh_{q,r}\bigr).
 \label{eq:fixed-limit}
\end{equation}

\begin{lemma}[Exact coefficient census]
\label{lem:coefficient-census}
Among the $512$ truth tables, exactly $192$ have $c_{11}=0$.  On those
tables,
\begin{align*}
 c_{01},c_{02},c_{12}&\in\{-1,0,1\},\\
 c_{21},c_{22}&\in\{-2,-1,0,1,2\}.
\end{align*}
There are $24$ distinct six-coefficient vectors, each with multiplicity
eight.  The table with plus-edge relation
\[
                       E=\{(0,-1),(0,+1)\}
\]
is self-compatible and has vector
$(c_{01},c_{02},c_{11},c_{12},c_{21},c_{22})=(1,0,0,0,-2,0)$.
\end{lemma}

\begin{proof}
This is exact ternary interpolation in the basis of
\eqref{eq:interpolation}.  The $c_{11}$ histogram over all $512$ tables is
\[
\begin{array}{c|rrrrr}
c_{11}&-1&-1/2&0&1/2&1\\ \hline
\text{count}&32&128&192&128&32.
\end{array}
\]
Exact row enumeration gives the alphabets and multiplicities.  For the
displayed relation, every plus edge has source $0$ and target nonzero, so
two plus edges cannot compose.  Direct interpolation gives the stated
vector.  All $4608$ point evaluations are reproduced independently in the
finite certificate; the enumeration is a finite identity, not asymptotic
evidence.
\end{proof}

\section{Squarefree cutoff and periodic Fourier bounds}

Fix an integer $P\ge2$ and put
\begin{equation}
 \eta_P(m)=\prod_{p\le P}\bigl(1-\mathbf1_{p^2\mid m}\bigr),
 \qquad
 M_P=\left(\prod_{p\le P}p\right)^2,
 \qquad
 \tau_P=\sum_{p>P}\frac1{p^2}.
 \label{eq:cutoff}
\end{equation}
Divisibility in \eqref{eq:cutoff} is defined for every integer.  In
particular,
\begin{equation}
                         \eta_P(0)=0,\qquad \eta_P(-1)=1.
 \label{eq:eta-endpoints}
\end{equation}
The function $\eta_P$ has period $M_P$.  This is a valid period; it need
not be the minimal period of a particular product involving $\eta_P$.

\begin{lemma}[Tail union bound]
\label{lem:tail}
For every positive integer $m$,
\begin{equation}
 0\le\eta_P(m)-\mob(m)^2
 \le\sum_{p>P}\mathbf1_{p^2\mid m}.
 \label{eq:point-tail}
\end{equation}
Consequently, for every $X\ge1$,
\begin{equation}
 \sum_{m\le X}\bigl(\eta_P(m)-\mob(m)^2\bigr)
 \le\sum_{p>P}\left\lfloor\frac X{p^2}\right\rfloor
 \le X\tau_P.
 \label{eq:summed-tail}
\end{equation}
The same inequalities compare every limiting periodic squarefree mean with
its $P$-cutoff mean.
\end{lemma}

\begin{proof}
If $\eta_P(m)=1$ but $\mob(m)^2=0$, a prime $p>P$ has $p^2\mid m$.
This proves \eqref{eq:point-tail}; summation gives
\eqref{eq:summed-tail}.  Notice that the exact count is
$\lfloor X/p^2\rfloor$, so no $\pi(\sqrt X)$ or $X^{-1/2}$ loss appears.
Divide the identical inequalities by $X$ and take limits along any fixed
periodic phase selection to obtain the last assertion.
\end{proof}

For a $Q$-periodic complex function $w$, use the normalized transform
\begin{equation}
 \widehat w(a)=\frac1Q\sum_{r\bmod Q}w(r)\e^{-2\pi i ar/Q}.
 \label{eq:dft}
\end{equation}

\begin{lemma}[Periodic M\"obius channel]
\label{lem:fourier}
Let
\[
 D(X)=\sup_{\alpha\in\mathbb R}
       \abs{\sum_{n\le X}\mob(n)\e^{2\pi i\alpha n}}.
\]
If $w$ has period $Q$, then
\begin{equation}
 \abs{\sum_{n\le X}w(n)\mob(n)}
 \le\sqrt Q\norm{w}_\infty D(X).
 \label{eq:fourier-bound}
\end{equation}
\end{lemma}

\begin{proof}
Fourier inversion and the definition of $D(X)$ give the right side with
$\sum_a\abs{\widehat w(a)}$ in place of
$\sqrt Q\norm{w}_\infty$.  Parseval and Cauchy--Schwarz give
\[
 \sum_a\abs{\widehat w(a)}
 \le\sqrt Q\left(\sum_a\abs{\widehat w(a)}^2\right)^{1/2}
 =\sqrt Q\left(\frac1Q\sum_r\abs{w(r)}^2\right)^{1/2}
 \le\sqrt Q\norm{w}_\infty.
\]
The sup norm is essential.  For $P=2$, $q=1$, and the legal
$c_{21}=-2$ witness, one period of the relevant shifted mask is
$(-2,-2,0,-2)$.  Its normalized Fourier $\ell^1$ norm is $3$, which is
larger than $\sqrt4=2$ but no larger than $2\sqrt4$.
\end{proof}

\section{The uniform analytic ledger}

For $N\ge3$, set
\begin{equation}
 D_*(N)=\max_{X\in\{N,N-2\}}D(X),
 \qquad Q=\operatorname{lcm}(q,M_P).
 \label{eq:D-star}
\end{equation}
Every coefficient sequence $c_{ij}(n\bmod q)$ and every cutoff mask in the
proof is $Q$-periodic.  Again, $Q$ is only asserted to be a common period.

\begin{proposition}[Exact conservative ledger]
\label{prop:ledger}
For every $N\ge3$, $P\ge2$, $q\ge1$, and $f\in\Fcal_q$,
\begin{equation}
 \boxed{
 \abs{S_N(q,f)-L_q(f)}
 \le\frac{4\sqrt QD_*(N)}N+13\tau_P+\frac{6Q}N+\frac4N.}
 \label{eq:master-bound}
\end{equation}
\end{proposition}

\begin{proof}
Insert \eqref{eq:interpolation} into \eqref{eq:finite-score}.  The
$c_{11}$ channel is absent phase by phase.  We account for the remaining
five channels separately.

The $c_{01}$ channel is a $Q$-periodic weight of sup norm one against
$\mob(n)$.  Lemma~\ref{lem:fourier} bounds it by
$\sqrt QD(N)/N$.

For $c_{12}$, the first two sites vanish and the change of variables
$m=n-2$ gives
\[
 \sum_{m\le N-2}c_{12}(m+2)\mob(m)\mob(m+2)^2.
\]
Replace $\mob(m+2)^2$ by $\eta_P(m+2)$.  The cost is at most
$N\tau_P$ by Lemma~\ref{lem:tail}.  The remaining periodic weight has sup
norm one, so its Fourier cost is $\sqrt QD(N-2)$.

For $c_{21}$, replace $\mob(n-2)^2$ by $\eta_P(n-2)$ for $n\ge3$.
The coefficient bound $\abs{c_{21}}\le2$ makes the tail cost
$2N\tau_P$.  Extend the periodic M\"obius sum to $1\le n\le N$.
At $n=2$, $\eta_P(0)=0$; at $n=1$, the cutoff has
$\eta_P(-1)=1$ whereas $\mob_0(-1)=0$, costing at most two.  The periodic
weight has sup norm two.  This channel therefore contributes
\[
                    2\sqrt QD(N)+2N\tau_P+2.
\]
The three Fourier costs are thus $1+1+2=4$.

It remains to compare the two square-only channels with
\eqref{eq:fixed-limit}.  For a bounded $Q$-periodic sequence $a$, complete
periods plus the final incomplete block give
\begin{equation}
 \abs{\frac1N\sum_{n\le N}a(n)-\frac1Q\sum_{r\bmod Q}a(r)}
 \le\frac{2Q\norm a_\infty}{N}.
 \label{eq:period-discrepancy}
\end{equation}
For $c_{02}\mob(n)^2$, finite-to-cutoff and
limit-to-cutoff each cost $\tau_P$, while
\eqref{eq:period-discrepancy} costs $2Q/N$.  Its total is
\begin{equation}
                         2\tau_P+2Q/N.
 \label{eq:c02-cost}
\end{equation}

For $c_{22}\mob_0(n-2)^2\mob(n)^2$, the product difference is bounded by
the sum of the two one-site differences.  Since
$\abs{c_{22}}\le2$, finite-to-cutoff costs $4\tau_P$ and
limit-to-cutoff costs another $4\tau_P$.  The cutoff product has sup norm
two, so \eqref{eq:period-discrepancy} costs $4Q/N$.  Extending the cutoff
product across the padded sites costs two at $n=1$ and zero at $n=2$ by
\eqref{eq:eta-endpoints}.  This channel contributes
\begin{equation}
                         8\tau_P+4Q/N+2/N.
 \label{eq:c22-cost}
\end{equation}

Collecting the tail coefficients gives
\[
 \underbrace{1+2}_{\text{Fourier masks}}
 +\underbrace{2}_{c_{02}}
 +\underbrace{8}_{c_{22}}=13.
\]
The periodic discrepancies give $2+4=6$, and the two endpoint costs give
$2+2=4$.  After division by $N$, this is
\eqref{eq:master-bound}.
\end{proof}

The decomposition also explains the tail constant in a second way.  The
finite mask replacements cost $1+1+2+4=8$ across the
$c_{12},c_{02},c_{21},c_{22}$ channels, and the two limiting square means
cost $1+4=5$.  These are conservative uniform costs; no sign cancellation
between channels is used.

\section{Polylogarithmic-clock uniformization}

All logarithms in the clock budget are natural.  For fixed $B>0$, put
\begin{equation}
                    H_B(N)=\left\lfloor(\log N)^B\right\rfloor.
 \label{eq:clock-budget}
\end{equation}
All suprema and maxima using this budget are understood for sufficiently
large $N$, when $H_B(N)\ge1$.

\begin{theorem}[Uniform polylogarithmic-clock limit]
\label{thm:uniform}
For every fixed real $B>0$,
\begin{equation}
 \boxed{
 \sup_{1\le q\le H_B(N)}\ \sup_{f\in\Fcal_q}
       \abs{S_N(q,f)-L_q(f)}\longrightarrow0.}
 \label{eq:uniform-limit}
\end{equation}
\end{theorem}

\begin{proof}
Take
\begin{equation}
                   P=P_N=\left\lfloor\sqrt{\log\log N}\right\rfloor,
 \label{eq:P-choice}
\end{equation}
which is at least two for all sufficiently large $N$.  The elementary
estimate
\begin{equation}
 \log M_P=2\sum_{p\le P}\log p\le2P\log P=o(\log\log N)
 \label{eq:M-growth}
\end{equation}
shows $M_P=(\log N)^{o(1)}$.  Uniformly for $q\le H_B(N)$,
\begin{equation}
               Q\le qM_P\le(\log N)^{B+o(1)}.
 \label{eq:Q-growth}
\end{equation}

RH-366 freezes Davenport's estimate: for every fixed $A>0$,
\begin{equation}
 \sup_{\alpha\in\mathbb R}
 \abs{\sum_{n\le X}\mob(n)\e^{2\pi i\alpha n}}
 \ll_A X(\log X)^{-A}.
 \label{eq:davenport}
\end{equation}
Choose one fixed $A>B/2$.  Equations \eqref{eq:Q-growth} and
\eqref{eq:davenport} make the Fourier term in
\eqref{eq:master-bound}
\[
                \ll_A(\log N)^{B/2-A+o(1)}=o(1).
\]
Also
\[
 \tau_P\le\sum_{n>P}n^{-2}\le\frac1{P-1}\longrightarrow0,
 \qquad Q/N\longrightarrow0,
 \qquad 1/N\longrightarrow0.
\]
The right side of \eqref{eq:master-bound} is independent of $f$ and
otherwise depends on $q$ only through the displayed uniform bound for $Q$.
This proves \eqref{eq:uniform-limit}.
\end{proof}

The fixed nature of $B$ is load-bearing.  The proof does not supply a
constant uniform in a varying $B=B(N)$.  For $q=N^\varepsilon$, the
$\sqrt Q$ Fourier mass is polynomial, while \eqref{eq:davenport} supplies
only an arbitrarily fixed logarithmic saving; this route therefore stops
before polynomial clocks.

\section{Triangular optimizers and a diagonal witness}

Because every $\Fcal_q$ is finite, define
\begin{equation}
 G_N(q)=\max_{f\in\Fcal_q}\abs{S_N(q,f)},
 \qquad
 G(q)=\max_{f\in\Fcal_q}\abs{L_q(f)}.
 \label{eq:optimizers}
\end{equation}
The second quantity is exactly the RH-379 fixed-clock optimizer.

\begin{corollary}[Uniform optimizer transfer]
\label{cor:optimizer-transfer}
For every fixed $B>0$,
\begin{equation}
 \sup_{1\le q\le H_B(N)}\abs{G_N(q)-G(q)}\longrightarrow0.
 \label{eq:optimizer-uniform}
\end{equation}
\end{corollary}

\begin{proof}
For any two finite families of real numbers $(a_f)$ and $(b_f)$,
\[
 \abs{\max_f\abs{a_f}-\max_f\abs{b_f}}
 \le\max_f\abs{a_f-b_f}.
\]
Apply this at each $q$ and then Theorem~\ref{thm:uniform}.
\end{proof}

\begin{theorem}[Restricted triangular-array endpoint]
\label{thm:triangular-endpoint}
For every fixed $B>0$,
\begin{equation}
          \boxed{\displaystyle
          \max_{1\le q\le H_B(N)}G_N(q)\longrightarrow B_\infty.}
 \label{eq:triangular-endpoint}
\end{equation}
\end{theorem}

\begin{proof}
Equation \eqref{eq:rh379-supremum} gives $G(q)\le B_\infty$ for every
fixed $q$.  Conversely, for every $\varepsilon>0$ there is one fixed clock
$q_\varepsilon$ with $G(q_\varepsilon)>B_\infty-\varepsilon$.
Since $H_B(N)\to\infty$, that clock is eventually admissible.  Hence
\[
              \max_{q\le H_B(N)}G(q)\longrightarrow B_\infty.
\]
Corollary~\ref{cor:optimizer-transfer} transfers this scalar maximum to
$G_N$.
\end{proof}

There is also an explicit cofinal witness, with a necessary empty-set
convention.  Let $3=p_1<p_2<\cdots$ be the odd primes and set
\begin{equation}
                       q_y=4\prod_{i\le y}p_i^2.
 \label{eq:square-clocks}
\end{equation}
When $H_B(N)<q_1=36$, record
\texttt{no\_square\_clock\_available}; do not silently substitute $q_1$.
Once the set is nonempty, define
\begin{equation}
 y_B(N)=\max\{y\ge1:q_y\le H_B(N)\}.
 \label{eq:y-diagonal}
\end{equation}
For each $y$, choose the RH-379 positive optimizer
$f_y^+\in\Fcal_{q_y}$ with $L_{q_y}(f_y^+)=G(q_y)>0$.

\begin{corollary}[Square-clock diagonal]
\label{cor:diagonal}
After \eqref{eq:y-diagonal} becomes defined,
\begin{equation}
 S_N\bigl(q_{y_B(N)},f_{y_B(N)}^+\bigr)\longrightarrow B_\infty.
 \label{eq:diagonal-limit}
\end{equation}
\end{corollary}

\begin{proof}
For every fixed $y$, eventually $q_y\le H_B(N)$; hence
$y_B(N)\to\infty$.  RH-379 proves $G(q_y)\to B_\infty$.  Therefore
\begin{align*}
 \abs{S_N(q_{y_B},f_{y_B}^+)-B_\infty}
 &\le\sup_{q\le H_B(N),\ f\in\Fcal_q}\abs{S_N(q,f)-L_q(f)}\\
 &\quad+\abs{G(q_{y_B})-B_\infty}\longrightarrow0.
\end{align*}
\end{proof}

The families $f_{y_B(N)}^+$ form a triangular-array witness only.  No
projective compatibility across clocks, and no single infinite selector,
is asserted.

\section{Finite certificate, source lock, and reproducibility}

The standard-library artifact reproduces the finite interfaces used above.
It contains all $512$ truth rows and all $4608$ interpolation evaluations;
the $192$ zero-$c_{11}$ rows reduce to $24$ coefficient vectors of
multiplicity eight.  It checks maximum coefficient-vector $\ell^1$ norm
$3$, maximum squared $\ell^2$ norm $5$, the self-compatible
$c_{21}=-2$ witness, and the three normalized-DFT channel costs $1,1,2$.

The cutoff rows have exact periods
\[
                         4,\quad36,\quad900,\quad44100
\]
for $P=2,3,5,7$.  Coprime and noncoprime least-common-multiple fixtures are
included.  One additional fixture has $q=5$, $P=3$, and $Q=180$, while its
relevant mask has minimal period $36$; this prevents the valid common
period from being promoted to a minimal-period claim.  Exact square-mask
means, the $13=8+5$ tail split, the $4$ endpoint cost, and small-clock
finite max-plus rows are regenerated from source.

The certificate is explicitly labelled
\texttt{reproduction\_not\_analytic\_proof}.  Its finite rows do not prove
Davenport's theorem, a uniform asymptotic, or any limit in this paper.  A
closed official Draft 2020--12 schema fixes every nested member and exact
JSON primitive type.  Twenty-four nontrivial certificate mutations test
coefficient, period, shift, factor-two, square-mean, tail, padding, clock,
sentinel, and claim-boundary failures.  Separate tests mutate source-lock
commits and group digests.  Optimized Python mode and bool-for-int
substitutions are also tested.

The source lock contains $67$ unique files in groups $51/8/8$: the released
RH-384 immutable closure, the standard eight RH-384 theorem files, and the
standard eight RH-366 Davenport files.  Every live file is compared with
its Git release blob.  Mutable root policy and handoff files are excluded.
The RH-366 release is
\texttt{0396fab97bbe3348c8237f8734dec0e1893fd3bf}.  Its frozen
\texttt{main.tex} SHA-256 is
\begin{center}
\footnotesize\ttfamily
7df165bd63d43f52dc217dea6691d231d8e40c00c148ab7e1aa4abcac55060fb
\end{center}
The RH-379 release is
\texttt{9ae9802ed17529ef4adfb81d7e2158d47c3c8d22}.  Its frozen
\texttt{main.tex} SHA-256 is
\begin{center}
\footnotesize\ttfamily
c5d97a227398a4f1d46a39fdec73ffb86aeb9bfc0f16296be7023b187b497090
\end{center}

\section{Context, limitations, and Gates}

The cubic residual limits recorded in RH-384 are useful context, but they
are not the discovery claim here \cite{RH384}.  The new theorem is the
restricted growing-clock uniformization in
Theorem~\ref{thm:uniform}.

Route A is \texttt{GO}.  Proposition~\ref{prop:ledger}, the fixed-$B$
cutoff, optimizer transfer, endpoint theorem, and diagonal witness form a
closed theorem chain.  Route B is \texttt{STOP\_SCOPED}.  The proof gives
no polynomial-clock uniformity: a fixed logarithmic Davenport saving does
not pay polynomial Fourier mass.  If $c_{11}(r)$ is activated, already
$q=1$ exposes the ordinary shift-two Chowla sum, which the frozen sources
do not cancel.

The following boundaries are mandatory.
\begin{enumerate}[leftmargin=2em]
 \item $B>0$ is fixed before $N\to\infty$; no $B=B(N)$ theorem is proved.
 \item No effective finite threshold or rate is claimed.  Constants in
       Davenport's estimate depend on the chosen fixed exponent.
 \item The maximum in \eqref{eq:triangular-endpoint} is restricted to
       $q\le H_B(N)$; it is not $\sup_{q<\infty}G_N(q)$.
 \item The theorem optimizes periodic phase tables, not arbitrary adaptive
       sign words.  It supplies no upper bound proving a limit for $K_N/N$.
 \item The diagonal witness is not a projectively compatible infinite
       selector.
\end{enumerate}

Gates A--E remain false/open.  These externally prescribed arithmetic
tables do not construct a canonical intrinsic dynamical spectral
determinant, time-oriented scattering completion, self-adjoint generator
with an intrinsic $T\log T$ law, signed von-Mangoldt weighted prime-power
trace, or completed-zeta divisor equality.  Nothing here identifies
Riemann zeros, constructs a Hilbert--P\'olya operator, or proves the Riemann
hypothesis.

\section{Conclusion}

The fixed-clock RH-379 limit is uniform throughout every fixed
polylogarithmic clock window.  The proof pays every source of growth: four
units of Fourier mass, thirteen units of squarefree-tail loss, six units of
period discrepancy, and four zero-padding endpoints.  The cutoff period is
small enough to fit inside Davenport's logarithmic saving precisely because
$B$ is fixed.  This establishes a genuine, restricted growing-clock theorem
while leaving polynomial clocks, active shift-two correlations, adaptive
capacity, and all spectral Gates open.

\section*{Reproducibility and declarations}

\textbf{Data availability.}  All source code, exact outputs, tests, closed
schema, source hashes, and archive-verification files are included in the
RH-385 repository directory.  No new external dataset was created or
analyzed.  Finite certificate rows reproduce interfaces and are not the
analytic proof.

\textbf{Ethics.}  This mathematical study uses no human participants,
animals, personal data, or biological materials.  Ethics approval was not
required.

\textbf{Author contributions.}  The RH research program performed
conceptualization, methodology, formal analysis, software, validation, and
writing---original draft and review/editing.

\textbf{Funding.}  This research received no specific grant from any
funding agency in the public, commercial, or not-for-profit sectors.

\textbf{Declaration of interests.}  The authors declare no known competing
financial interests or personal relationships that could have influenced
the work reported here.

\textbf{AI-assisted workflow.}  The manuscript and executable checks were
prepared with AI-assisted academic writing and coding tools under the
repository's documented multi-agent workflow.  The mathematical claims are
bounded by frozen sources, explicit proofs, independent audits, and
reproducible artifacts.  The authors retain responsibility for the final
content and its integrity.

\begingroup
\small
\bibliographystyle{plain}
\bibliography{references}
\endgroup
\end{document}
